\chapter{Integrated Data Repository Architecture}\label{appendix:architecture} 
\section{Repository Layers} 

The Integrated Data Repository (IDR) adopted a layered architecture consisting of source systems, staging tables, intermediate integration views, analytical fact tables, and reporting dashboards. This design separated raw operational data from analytical representations while maintaining data provenance and traceability. 

\section{Primary Repository Components} 

\begin{table}[h] 
    \centering 
    \begin{tabular}{|p{4cm}|p{8cm}|} 
        \hline 
        \textbf{Component} & \textbf{Purpose} \\ 
        \hline 
        Staging Layer & Preservation of raw source extracts and audit trail. \\ 
        \hline 
        Integration Layer & Application of business rules and data harmonisation processes. \\
        \hline 
        Fact Tables & Analytics-ready datasets used for operational reporting. \\ 
        \hline 
        Power BI Layer & Operational dashboards and data visualisation. \\ 
        \hline 
    \end{tabular} 
    \caption{Repository Architecture Components} 
\end{table} 

\section{Source Systems} 

\begin{table}[h] 
    \centering 
    \begin{tabular}{|p{3cm}|p{4cm}|p{5cm}|} 
        \hline 
        \textbf{Source} & \textbf{Data Type} & \textbf{Purpose} \\ 
        \hline 
        Oncology Referral Register & SharePoint List & Oncology referral management \\ 
        \hline Patient Administration System (PAS) & Operational Reports & Clinic bookings and attendance \\ 
        \hline 
        Aria OIS & CSV Extracts & Radiotherapy referral, CT and treatment activity \\ 
        \hline 
    \end{tabular} 
    \caption{Primary Data Sources} 
\end{table} 

\section{Logical Architecture} 

Figure \ref{fig:LogicalArch} in Chapter \ref{method} illustrates the logical architecture of the repository.